\section{Introduction}
\label{sec:intro}

% Outline (from docs/paper-plan.md "Why this paper, why now"):
% 1. SmaTable: hidden piezo sensors in a table surface, privacy-by-design gesture input.
% 2. Deployment exposes three continual-learning questions the literature does not answer
%    for this system: user adaptation budget, new-gesture budget, replay budget.
% 3. Hettstedt et al. 2026 established the DAQ chain on the RP2350 but inference and
%    training stayed on a PC; Yoshida 2023a/b are PC-class and need calibration data to
%    leave the device.
% 4. This paper: first on-RP2350 inference + training measurements, and quantitative
%    sample-budget / stability envelopes for three CL scenarios.

\todo{Motivation paragraph: SmaTable concept and why on-device training matters for it (privacy: calibration data must not leave the table).}

\todo{Gap paragraph: prior SmaTable work \cite{hettstedt2026smatable,yoshida2023a,yoshida2023b} keeps training on a PC; no on-device numbers exist.}

\paragraph{Contributions.}
\begin{enumerate}
  \item \textbf{Replay sizing curve (RQ1).} Accuracy on the original users as a function of per-class replay memory under three replay mechanisms at matched bytes per class --- random exemplars, herding exemplars~\cite{rebuffi2017icarl}, and PPCA generative replay~\cite{tipping1999ppca} --- plus a no-replay floor, measured on-device on the RP2350.
  \item \textbf{User-adaptation sample curve (RQ2).} Accuracy on an unseen user as a function of the number of that user's samples used for on-device fine-tuning, under leave-one-subject-out cross-validation, reporting target-user gain, original-user retention and forward transfer~\cite{lopezpaz2017gem}.
  \item \textbf{New-gesture sample curve (RQ3).} Accuracy on one added gesture class as a function of samples per contributor $\times$ number of contributors, while bounding degradation on the pre-existing classes.
  \item \textbf{On-device training stability (RQ4).} Run-to-run variance of final accuracy across seeds, decomposed into new-parameter initialisation and on-device data order.
  \item \textbf{Augmentation effect (RQ5).} Whether additive Gaussian noise reduces the sample budgets of RQ1/RQ2 at matched accuracy.
\end{enumerate}

\todo{One sentence on what is explicitly \emph{not} claimed: framework additions (Conv1d, LayerNorm, GroupNorm, PPCA replay in the ODT library) are enablers; architecture is fixed; no QPU.}
